% ===========================================================================
% Poste 1 — RRQ
% ===========================================================================
\poste{1}{Régime de rentes du Québec (RRQ)}{CA\_rrq}

\pdesc Cotisation obligatoire prélevée sur le revenu de travail, qui finance
les rentes de retraite, d'invalidité et de survivant du régime public
québécois. Depuis 2019, le régime comporte un volet \emph{de base} et un
volet \emph{supplémentaire} (la \og bonification \fg{}), dont le traitement
fiscal diffère.

\pobj Assurer un revenu de remplacement à la retraite (et en cas
d'invalidité ou de décès), proportionnel aux gains de travail.

\pregle La cotisation d'un adulte se calcule sur deux bandes de son revenu de
travail $r$. La première bande court de l'exemption générale $E$ au maximum
des gains admissibles ($MGA$) ; la seconde, du $MGA$ au maximum
supplémentaire ($MGAS$, fixé à \pc{114} du $MGA$) :
\begin{align*}
B_1 &= \pos{\min(r,\, MGA) - E} \\
B_2 &= \pos{\min(r,\, MGAS) - MGA}
\end{align*}
Trois taux s'appliquent (part de l'employé) : le taux du régime de base
$\tau_{base}$ et la première cotisation supplémentaire $\tau_1$ sur la
bande~1 ; la deuxième cotisation supplémentaire $\tau_2$ sur la bande~2 :
\begin{align*}
\text{cotisation de base} &= \tau_{base} \times B_1 \\
\text{cotisation supplémentaire} &= \tau_1 \times B_1 + \tau_2 \times B_2 \\
\text{cotisation totale} &= \text{base} + \text{supplémentaire}
\end{align*}
La cotisation du ménage est la somme des cotisations des adultes. Les
ménages retraités ne cotisent pas (aucun revenu de travail).

La distinction base/supplémentaire importe pour l'impôt (poste~19) : la
cotisation \emph{de base} donne droit à un crédit d'impôt non remboursable au
fédéral, tandis que la cotisation \emph{supplémentaire} est déductible du
revenu --- elle réduit le revenu imposable et le revenu net aux deux ordres
de gouvernement.

\emph{Point de contrôle.} En 2025, un salarié gagnant le $MGAS$ ou plus paie
au maximum : base $= 0{,}054 \times (71\,300 - 3\,500) = \D{3661.20}$ ;
supplémentaire $= 0{,}01 \times 67\,800 + 0{,}04 \times 9\,900 = \D{1074}$ ;
total $= \D{4735.20}$.

\pparams

\begin{center}
\small
\begin{tabular}{l r r}
\toprule
Paramètre & 2025 & 2026 \\
\midrule
Exemption générale ($E$)                          & \D{3500}  & \D{3500} \\
Maximum des gains admissibles ($MGA$)             & \D{71300} & \D{74600} \\
Maximum supplémentaire ($MGAS$)                   & \D{81200} & \D{85000} \\
Taux du régime de base ($\tau_{base}$)            & \pc{5.4}  & \pc{5.3} \\
Première cotisation supplémentaire ($\tau_1$)     & \pc{1}    & \pc{1} \\
Deuxième cotisation supplémentaire ($\tau_2$)     & \pc{4}    & \pc{4} \\
\bottomrule
\end{tabular}
\end{center}

La baisse du taux de base de \pc{5.4} à \pc{5.3} en 2026 découle de
l'évaluation actuarielle du régime ; c'est l'un des allègements visibles dans
l'écart 2025--2026 du calculateur.

% ============================================================================
% GÉNÉRÉ par scripts/exemples-guide.ts — NE PAS ÉDITER À LA MAIN.
% Régénérer : npm run exemples (chaque chiffre est vérifié contre le moteur).
% ============================================================================
\pexemples La cotisation se calcule \emph{par adulte}, sur le revenu de travail.
Paramètres 2025 : exemption de \D{3500}, MGA de \D{71300}, MGAS de
\D{81200}, taux de \pc{5.4} (base), \pc{1} (1\iere{}
supplémentaire) et \pc{4} (2\ieme{} supplémentaire).

\medskip\noindent\textbf{M3 --- personne vivant seule, 30~ans, revenu de travail de \D{15000}.}
Le revenu est sous le MGA : seule la bande~1 est occupée.
\begin{align*}
\text{bande 1} &= \min(\num{15000},\ \num{71300}) - \num{3500} = \num{11500} \\
\text{régime de base} &= \num{11500} \times \pc{5.4} = \num{621} \\
\text{1\iere{} cotisation supplémentaire} &= \num{11500} \times \pc{1} = \num{115} \\
\text{bande 2} &= 0 \quad (\text{revenu} \le \text{MGA}) \\
\text{cotisation totale} &= \num{736}
\end{align*}
Cotisation RRQ : \D{736}, dont \D{621} de régime de base
(créditable à l'impôt) et \D{115} de régime supplémentaire (déductible).

\medskip\noindent\textbf{M4 --- personne vivant seule, 40~ans, revenu de travail de \D{50000}.}
Même mécanique, plus haut dans la bande~1.
\begin{align*}
\text{bande 1} &= \min(\num{50000},\ \num{71300}) - \num{3500} = \num{46500} \\
\text{régime de base} &= \num{46500} \times \pc{5.4} = \num{2511} \\
\text{1\iere{} cotisation supplémentaire} &= \num{46500} \times \pc{1} = \num{465} \\
\text{bande 2} &= 0 \quad (\text{revenu} \le \text{MGA}) \\
\text{cotisation totale} &= \num{2976}
\end{align*}
Cotisation RRQ : \D{2976} (base \D{2511} ; supplémentaire \D{465}).

\medskip\noindent\textbf{M5 --- personne vivant seule, 45~ans, revenu de travail de \D{100000}.}
Le revenu dépasse le MGAS (\D{81200}) : les deux bandes sont
\emph{saturées} — la cotisation atteint son maximum 2025.
\begin{align*}
\text{bande 1} &= \min(\num{100000},\ \num{71300}) - \num{3500} = \num{67800} \\
\text{régime de base} &= \num{67800} \times \pc{5.4} = \num{3661.20} \\
\text{1\iere{} cotisation supplémentaire} &= \num{67800} \times \pc{1} = \num{678} \\
\text{bande 2} &= \min(\num{100000},\ \num{81200}) - \num{71300} = \num{9900} \\
\text{2\ieme{} cotisation supplémentaire} &= \num{9900} \times \pc{4} = \num{396} \\
\text{cotisation totale} &= \num{4735.20}
\end{align*}
Cotisation RRQ maximale : \D{4735.20}, dont \D{1074} de régime
supplémentaire (déductible du revenu imposable, ce qui réduit aussi les bases
$\RFN$ et $\AFNI$).

\medskip\noindent\textbf{M8 --- couple, 38 et 36~ans, revenus de travail de \D{60000} et \D{40000}, deux enfants : 4~ans (garde non subventionnée, \D{13000}) et 8~ans (garde non subventionnée, \D{3000}).}
Pour un couple, on calcule chaque adulte séparément puis on additionne :
\begin{align*}
\text{adulte 1 } (\num{60000}) &: \num{3616} \\
\text{adulte 2 } (\num{40000}) &: \num{2336} \\
\text{ménage} &= \num{3616} + \num{2336} = \num{5952}
\end{align*}
Cotisation RRQ du ménage : \D{5952}.


\prefs Loi sur le régime de rentes du Québec (RLRQ, c.~R-9) ; Retraite
Québec (taux et maximums) ; Revenu Québec.
